\section{Naive Exploration}

\begin{frame}{Greedy Algorithm}
\begin{itemize}
    \item For each arm, maintain a \textbf{sample-mean estimate} of its value --- average the rewards seen from that arm so far:
    \[
        \hat{Q}_t(a) \;\stackrel{\text{def}}{=}\; \frac{1}{N_t(a)} \sum_{s=1}^{t} r_s \cdot \mathbf{1}_{a_s = a} \;\;\approx\;\; Q(a)
    \]
    \item Pick the action with highest estimate:
    \[
        a_t = \operatorname*{arg\,max}_{a \in \mathcal{A}} \hat{Q}_{t-1}(a)
    \]
    \item \textcolor{red}{Greedy can lock onto a suboptimal action forever $\Rightarrow$ \textbf{linear total regret}}
\end{itemize}
\end{frame}

\begin{frame}{$\epsilon$-Greedy Algorithm}
\begin{itemize}
    \setlength{\itemsep}{0.5em}
    \item Force a small amount of exploration at every step. At time $t$:
    \begin{itemize}
        \item With probability $1 - \epsilon$, select $a_t = \operatorname*{arg\,max}_{a} \hat{Q}_{t-1}(a)$
        \item With probability $\epsilon$, select a uniformly random action
    \end{itemize}
    \item Constant $\epsilon$ ensures a minimum per-step regret proportional to the mean gap:
    \[
        l_t \;\geq\; \frac{\epsilon}{|\mathcal{A}|} \sum_{a \in \mathcal{A}} \Delta_a
    \]
    \item \textcolor{red}{Therefore $\epsilon$-Greedy also has \textbf{linear total regret}}
\end{itemize}
\end{frame}

% Day-2 callback as its own small slide (mirrors the Day-3 REINFORCE callback)
\begin{frame}{$\epsilon$-Greedy Is DQN's Exploration Rule}
\vspace{0.5em}
\begin{center}
\textcolor{red}{\textbf{This is the exact exploration rule used in DQN on Day 2.}}
\end{center}
\vspace{1em}
\begin{itemize}
    \setlength{\itemsep}{0.6em}
    \item DQN inherits $\epsilon$-Greedy's \emph{linear-regret} weakness
    \item It works in practice with a decay schedule (next slide), but it has no theoretical guarantee of efficient exploration
    \item This is why the rest of this lecture exists --- everything that follows is a more principled answer to ``how do I explore?''
\end{itemize}
\end{frame}

\begin{frame}{Optimistic Initialization}
\begin{itemize}
    \setlength{\itemsep}{0.45em}
    \item Simple, practical idea: initialize $\hat{Q}_0(a)$ to a \emph{high} value for every $a$
    \item Update by incremental averaging:
    \[
        \hat{Q}_t(a_t) = \hat{Q}_{t-1}(a_t) + \frac{1}{N_t(a_t)} \big(r_t - \hat{Q}_{t-1}(a_t)\big)
    \]
    \item Encourages systematic exploration early on (every arm starts off looking great until tried)
    \item But the algorithm can still lock onto a suboptimal action $\Rightarrow$ \textbf{linear regret}
\end{itemize}
\end{frame}

\begin{frame}{Decaying $\epsilon_t$-Greedy}
\begin{itemize}
    \setlength{\itemsep}{0.65em}
    \item Idea: let $\epsilon$ shrink over time --- explore aggressively early, exploit increasingly late
    \item In theory, a schedule that depends on the (unknown) suboptimality gaps $\Delta_a$ achieves \emph{logarithmic} regret --- but you can't compute it without knowing the gaps, so it's impractical
    \item \textcolor{red}{In practice} use problem-agnostic schedules: $\epsilon_t = 1/t$, $\epsilon_t = 1/\sqrt{t}$, or simple linear annealing --- no guarantee, but all work in practice
    \item \href{https://github.com/TikhonJelvis/RL-book/blob/master/rl/chapter14/epsilon_greedy.py}{Educational Code} for decaying $\epsilon$-greedy with optimistic initialization
\end{itemize}
\end{frame}

\begin{frame}{The Lai-Robbins Lower Bound: Setup}
\textbf{What is a ``lower bound''?} A statement that \emph{no algorithm, however clever}, can do better than a certain regret growth rate --- it tells us the best we can ever hope for.
\vspace{0.5em}
\begin{itemize}
    \setlength{\itemsep}{0.55em}
    \item \textbf{Goal:} an algorithm with \emph{sublinear} total regret $L_T = \sum_{t=1}^T l_t$ for any MAB, without knowing the reward distributions $\mathcal{R}^a$ in advance
    \item Two quantities will appear:
    \begin{itemize}
        \item \textbf{Gap} $\Delta_a = V^* - Q(a)$: how much value we lose \emph{each time} we pull suboptimal arm $a$
        \item \textbf{KL-divergence} $\mathrm{KL}(\mathcal{R}^a \,\|\, \mathcal{R}^{a^*})$: a statistical distance between two distributions --- \emph{how easy is it to tell arm $a$ apart from the best arm $a^*$ from samples?} \\
              Small KL $\Rightarrow$ distributions look almost identical $\Rightarrow$ need many samples to distinguish
    \end{itemize}
\end{itemize}
\end{frame}

\begin{frame}{The Lai-Robbins Lower Bound: The Theorem}
\begin{block}{Theorem (Lai \& Robbins, 1985)}
For any reasonable\footnote{\scriptsize ``consistent'': sub-polynomial regret on every problem instance} bandit algorithm, asymptotically:
\[
    \lim_{T \to \infty} \frac{L_T}{\log T} \;\;\geq\;\; \sum_{a \mid \Delta_a > 0} \frac{\Delta_a}{\mathrm{KL}(\mathcal{R}^a \,\|\, \mathcal{R}^{a^*})}
\]
\end{block}
\begin{itemize}
    \setlength{\itemsep}{0.55em}
    \item \textbf{Reading it.} Each suboptimal arm contributes regret proportional to its \emph{gap}, scaled \emph{inversely} by how easy it is to tell apart from the best
    \item \textbf{The worst kind of arm:} large gap (expensive to pull) \emph{and} small KL (hard to distinguish, so we keep pulling it) --- exactly what you'd expect
    \item \textcolor{red}{The best regret growth any algorithm can achieve is $\mathcal{O}(\log T)$} --- UCB (next) and Thompson sampling both match this rate
\end{itemize}
\end{frame}
